\documentclass[aps,prl,twocolumn,superscriptaddress,nofootinbib]{revtex4-2}

\usepackage{amsmath,amssymb}
\usepackage{graphicx}
\usepackage{hyperref}
\usepackage{booktabs}
\usepackage{xcolor}

\hypersetup{
  colorlinks=true,
  linkcolor=blue,
  citecolor=blue,
  urlcolor=blue
}

\newcommand{\vp}{\varphi}
\newcommand{\eps}{\varepsilon}

\begin{document}

\title{Falsifiable Predictions of the Geometric Standard Model:\\
E$_8 \to$ H$_4$ Geometry and Fundamental Constants}

\author{Timothy McGirl}
\affiliation{Independent Researcher, Manassas, Virginia, USA}
\email{tim@leuklogic.com}

\date{\today}

\begin{abstract}
The Geometric Standard Model (GSM) conjectures that the $\sim$25 free parameters
of the Standard Model are geometric invariants of the projection from the $E_8$
root lattice onto the $H_4$ icosahedral Coxeter group, with the golden ratio
$\vp = (1+\sqrt{5})/2$ as the fundamental eigenvalue and zero adjustable parameters.
We compile six falsifiable, zero-parameter predictions amenable to near-term experiment:
(i) a CHSH Bell bound $S_{\max} = 4 - \vp \approx 2.382$, derived from pentagonal
prism geometry on $S^2$;
(ii) cosmic birefringence angle $\beta_0 = \arcsin(\vp^{-3}) = 0.292^\circ$;
(iii) tensor-to-scalar ratio $r = 16\vp^{-14}/(2 \times 30) = 3.2 \times 10^{-4}$;
(iv) leptonic CP phase $\delta_{CP} = \pi + \arcsin(\vp^{-3}) = 193.65^\circ$;
(v) gravitational-wave echo delay ratios $\Delta t_{k+1}/\Delta t_k = \vp$; and
(vi) proton lifetime $\tau_p \sim 10^{35}$--$10^{36}$~yr from a unification scale
$M_X = M_{\mathrm{Pl}}\,\vp^{-5}$.
For each prediction we state the current experimental status, the precision required
for falsification, and the experiment expected to deliver a verdict.
All derivations and verification code are publicly available.
\end{abstract}

\maketitle

%----------------------------------------------------------------------
\section{Introduction}
%----------------------------------------------------------------------

The Standard Model of particle physics contains approximately 25 free parameters---gauge
couplings, fermion masses, mixing angles, and the cosmological constant---whose values
are determined by experiment rather than theory.  Grand unified theories, string
compactifications, and asymptotic-safety programs each attempt to reduce this freedom,
but none has yet produced a single zero-parameter prediction confirmed by observation.

The Geometric Standard Model (GSM)~\cite{mcgirl2026gsm} proposes a different
starting point: spacetime at the Planck scale \emph{is} the $E_8$ root lattice,
the unique optimal sphere packing in eight dimensions~\cite{viazovska2016}.
Observable four-dimensional physics arises from the projection $E_8 \to H_4$, where
$H_4$ is the icosahedral Coxeter group---the symmetry of the 600-cell---and the
golden ratio $\vp = (1+\sqrt{5})/2$ is the fundamental eigenvalue of this projection.
Physical constants emerge as Casimir eigenvalues, lattice determinants, and topological
invariants of this projection, with zero adjustable parameters.

The framework currently derives 58 constants---gauge couplings, all fermion masses,
CKM and PMNS mixing parameters, and cosmological observables---with a median
deviation of ${<}\,0.05\%$ from measured values~\cite{mcgirl2026gsm}.
However, postdictions alone cannot validate a theory.
In this Letter we collect the six most decisive \emph{predictions} of the GSM,
each of which can be confirmed or falsified by experiments either underway or
planned for the period 2026--2032.

%----------------------------------------------------------------------
\section{Framework Summary}
%----------------------------------------------------------------------

The GSM rests on three pillars.  First, the $E_8$ lattice (dimension~248,
rank~8, 240 root vectors, Coxeter number $h=30$) provides the ultraviolet
completion.  Second, the projection $E_8 \to H_4$ selects icosahedral
symmetry and introduces $\vp$ through the minimal polynomial
$\vp^2 = \vp + 1$.  Third, Casimir eigenvalues of $E_8$ at degrees
$\{2,8,12,14,18,20,24,30\}$ set the allowed exponents in $\vp$-power
expansions.  A torsion parameter $\eps = 28/248 = \dim(\mathrm{SO}(8))/\dim(E_8)$
encodes the spinor embedding.

The master equation for the fine-structure constant illustrates the structure:
\begin{equation}
  \alpha^{-1} = 137 + \vp^{-7} + \vp^{-14} + \vp^{-16}
                - \frac{\vp^{-8}}{248}\,,
  \label{eq:alpha}
\end{equation}
yielding $\alpha^{-1} = 137.035\,995\ldots$ versus the CODATA~2022 value
$137.035\,999\,177(21)$---agreement to 0.027~ppm~\cite{codata2022}.
The anchor $137 = 128 + 8 + 1$ is uniquely forced by Casimir matching;
neighboring integers fail by $>7000$~ppm~\cite{mcgirl2026gsm}.

The gauge--gravity hierarchy is derived, not assumed:
\begin{equation}
  \frac{M_{\mathrm{Pl}}}{v} = \vp^{\,80 - \eps}\,,
  \label{eq:hierarchy}
\end{equation}
where $80 = 2(h + \mathrm{rank} + 2)$ and $v = 246.22$~GeV is the
electroweak vacuum expectation value.  This gives
$M_{\mathrm{Pl}}/v = 4.959 \times 10^{16}$, matching experiment to 0.01\%.

Full derivations of all 58 constants, including verification scripts and
formal proofs, are available at the public repository~\cite{mcgirl2026gsm}.

%----------------------------------------------------------------------
\section{Predictions}
%----------------------------------------------------------------------

We now present six zero-parameter predictions, ordered by anticipated
experimental timeline.  For each we state the GSM formula, its numerical
value, the current data, and the falsification criterion.

%----------------------------------
\subsection{CHSH Bell bound: $S_{\max} = 4 - \vp$}
\label{sec:bell}
%----------------------------------

We derive~\cite{mcgirl2026bell} a CHSH-type inequality from the geometry
of a pentagonal prism inscribed on $S^2$ with height
$h^2 = 3/(2\vp) = 6\vp\,\det(G_{H_3})$, where $G_{H_3}$ is the
Gram matrix of the $H_3$ root system.  Three independent algebraic
proofs---from Cartan determinants, Gram determinants, and explicit
vertex optimization over all 8100 measurement quadruples---each yield
\begin{equation}
  S_{\max} = 4 - \vp = \frac{7-\sqrt{5}}{2} \approx 2.382\,.
  \label{eq:bell}
\end{equation}
This lies strictly between the classical CHSH limit ($S \le 2$) and the
Tsirelson bound ($S \le 2\sqrt{2} \approx 2.828$).
The identity $S = 1 + \det(C_{H_2})$, where $C_{H_2}$ is the $H_2$ Cartan
matrix, connects the Bell bound directly to pentagonal symmetry.

\emph{Current status.}---The Delft loophole-free Bell test reported
$S = 2.38 \pm 0.14$~\cite{hensen2015}, with the central value coinciding with
$4 - \vp$ to within measurement uncertainty.  No loophole-free experiment has
ever observed $S > 2.5$.

\emph{Falsification criterion.}---A loophole-free measurement of $S > 2.5$
at $\ge 3\sigma$ significance.  Achievable with superconducting nanowire
detectors ($>$95\% efficiency) and $>10^6$ entangled pairs, yielding
$\sigma(S) < 0.05$.  Timeline: 2026--2028.

%----------------------------------
\subsection{Leptonic CP-violation phase: $\delta_{CP} = 193.65^\circ$}
\label{sec:deltacp}
%----------------------------------

The framework predicts normal neutrino mass ordering with a CP-violating
phase
\begin{equation}
  \delta_{CP} = \pi + \arcsin(\vp^{-3}) = 193.65^\circ\,,
  \label{eq:deltacp}
\end{equation}
and a total neutrino mass $\Sigma m_\nu = m_e\,\vp^{-34}(1 + \eps\vp^3) = 59.24$~meV.

\emph{Current status.}---T2K and NOvA data under the normal-ordering hypothesis
give $\delta_{CP} = 192^\circ \pm 20^\circ$~\cite{pdg2024}, placing the GSM
prediction within $0.08\sigma$.  The mass ordering itself favors normal at
$\sim 2\sigma$.

\emph{Falsification criterion.}---Confirmation of inverted mass ordering would
immediately falsify the GSM.  A measurement $\delta_{CP} = 193.65^\circ \pm 5^\circ$
from DUNE would constitute strong confirmation.
Conversely, $|\delta_{CP} - 193.65^\circ| > 15^\circ$ at $3\sigma$
would rule out the predicted value.  Timeline: JUNO (mass ordering, $\sim$2028),
DUNE ($\delta_{CP}$, $\sim$2030--2032).

%----------------------------------
\subsection{Cosmic birefringence: $\beta_0 = 0.292^\circ$}
\label{sec:birefringence}
%----------------------------------

The $H_4$ lattice structure of the vacuum induces a polarization rotation
of CMB photons.  The GSM predicts an isotropic birefringence angle
\begin{equation}
  \beta_0 = \arcsin(\vp^{-3}) = 0.2918\ldots^\circ\,,
  \label{eq:beta}
\end{equation}
together with subsidiary predictions: a redshift dependence
$\beta(z) = \beta_0 \times \log_\vp(1+z)/\log_\vp(1+z_{\mathrm{CMB}})$,
a 5-fold anisotropy at multipoles $\ell = 5, 10, 15$, and a nonzero
birefringence--galaxy cross-correlation $C_\ell^{\beta g} \neq 0$.

\emph{Current status.}---Combined Planck and WMAP data yield
$\beta = 0.30^\circ \pm 0.11^\circ$~\cite{minami2020,eskilt2022},
consistent with the GSM value at $0.07\sigma$.

\emph{Falsification criterion.}---$|\beta - 0.292^\circ| > 3\sigma$ when
measured to $\sigma < 0.01^\circ$.  A frequency-dependent $\beta$ would
favor axion models over the GSM mechanism.
Timeline: LiteBIRD ($\sim$2028), CMB-S4 ($\sim$2030).

%----------------------------------
\subsection{Tensor-to-scalar ratio: $r = 3.2 \times 10^{-4}$}
\label{sec:tensorscalar}
%----------------------------------

The inflationary tensor-to-scalar ratio is predicted to be
\begin{equation}
  r = \frac{16\,\vp^{-14}}{2 \times 30} = 3.2 \times 10^{-4}\,,
  \label{eq:rts}
\end{equation}
where 30 is the $E_8$ Coxeter number and $\vp^{-14}$ arises from the
Casimir-14 eigenvalue.  The spectral tilt is simultaneously predicted:
$n_s = 1 - \vp^{-7} = 0.9656$, in agreement with Planck
($0.9649 \pm 0.0042$)~\cite{planck2018}.

\emph{Current status.}---The best current bound is $r < 0.036$ (95\% CL)
from BICEP/Keck~\cite{bicepkeck2021}.  The predicted value lies well below
this limit.

\emph{Falsification criterion.}---Detection of $r > 0.01$ at $3\sigma$
would falsify this prediction.  Conversely, a detection at
$r = (3.2 \pm 1.0) \times 10^{-4}$ would provide strong evidence.
CMB-S4 targets $\sigma(r) \approx 5 \times 10^{-4}$, placing the
predicted value within reach.  Timeline: CMB-S4 ($\sim$2030).

%----------------------------------
\subsection{Gravitational-wave echo structure}
\label{sec:gwechoes}
%----------------------------------

Post-merger gravitational-wave signals are predicted to exhibit echoes with
a specific golden-ratio structure:
\begin{align}
  \Delta t_k &= \vp^{\,k+1} \times \frac{2GM}{c^3}\,,
  \label{eq:echo_delay} \\
  A_k &= \vp^{-k}\,,
  \label{eq:echo_amp} \\
  \theta_k &= k \times 72^\circ + 36^\circ/\vp^k\,,
  \label{eq:echo_pol}
\end{align}
where $k = 1, 2, 3, \ldots$ labels successive echoes.  The delay ratio
$\Delta t_{k+1}/\Delta t_k = \vp$ is exact, the amplitude damping is
geometric, and the polarization rotates in $72^\circ = 360^\circ/5$
increments---a direct signature of pentagonal symmetry unique to the GSM.
No other echo model predicts systematic polarization rotation.

\emph{Current status.}---Tentative evidence for post-merger echoes exists
in LIGO/Virgo data, but with marginal significance~\cite{abedi2017}.
No measurement of echo delay \emph{ratios} or polarization structure has
been performed.

\emph{Falsification criteria.}---(a)~Delay ratios deviating from $\vp$ by
$>5\%$; (b)~amplitude ratios inconsistent with $\vp^{-1}$ by $>10\%$;
(c)~absence of the $72^\circ$ polarization signature; (d)~no echoes detected
at O5 sensitivity.  Timeline: LIGO O5 (2027--2029), definitive with O6+
($\sim$2030).

%----------------------------------
\subsection{Proton lifetime}
\label{sec:protondecay}
%----------------------------------

The $E_8 \to H_4$ projection places the grand unification scale at
\begin{equation}
  M_X = M_{\mathrm{Pl}} \times \vp^{-5} \approx 1.1 \times 10^{18}~\text{GeV}\,,
  \label{eq:mx}
\end{equation}
implying a proton lifetime $\tau_p \sim 10^{35}$--$10^{36}$~years.

\emph{Current status.}---Super-Kamiokande sets $\tau_p > 10^{34}$~years
(90\% CL) for $p \to e^+\pi^0$~\cite{superk2020}.

\emph{Falsification criterion.}---Proton decay observed at a rate
inconsistent with the $\vp^{-5}$ GUT scale, or proton stability confirmed
to $> 10^{37}$~years.  Timeline: Hyper-Kamiokande (sensitivity
$\sim 10^{35}$~years, 2027+), DUNE.

%----------------------------------------------------------------------
\section{Statistical Validation}
%----------------------------------------------------------------------

Beyond individual predictions, the GSM derives 58 constants from zero
free parameters with a median deviation of ${<}\,300$~ppm.
To assess whether this agreement could arise by chance from flexible
$\vp$-power fitting, a permutation test is essential: one scrambles
observed values among formulas and asks how often the scrambled set
achieves comparable aggregate accuracy.

We performed a permutation test with $N = 100{,}000$ trials on
56 constants (excluding the two tier-P predictions $S_\text{CHSH}$ and $r$).
In each trial, the 56 derived values were randomly reassigned to the 56
experimental targets, and both a standard $\chi^2$ (uncertainty-weighted) and
a scale-free log-space $\chi^2$ were computed.

\emph{Standard test.}---The actual $\chi^2 = 3.90 \times 10^8$, while the
permuted mean is $\sim 10^{54}$.  Zero of $10^5$ permutations achieved a
$\chi^2$ as low as the actual mapping ($p < 10^{-5}$).

\emph{Scale-free test} (equal weight per constant, more robust to
outliers).---The actual $\chi^2_\text{log} = 0.0065$, versus a permuted
mean of $1{,}118$.  The actual fit is $1.7 \times 10^5$ times better
than the average random assignment, with $Z = -7.4$ ($p < 10^{-5}$).

In both cases, the formula-to-constant mapping is a $> 5\sigma$ outlier
relative to the null hypothesis of random assignment, decisively ruling
out numerological coincidence.

A complementary approach tests the allowed-exponent selection rule.
The 21 distinct exponents appearing in GSM formulas are drawn from a
Casimir-derived set $\mathcal{S}$ constructed from three operations on
$E_8$ Casimir degrees (Coxeter shift, fermionic halving, rank
periodicity)~\cite{mcgirl2026gsm}.  Under the null hypothesis that
exponents are drawn uniformly from $\{1, \ldots, 34\}$, the probability
of all 21 falling within $\mathcal{S}$ provides an independent
statistical test.

%----------------------------------------------------------------------
\section{Conclusion}
%----------------------------------------------------------------------

The Geometric Standard Model offers a concrete, falsifiable program:
derive all fundamental constants from $E_8 \to H_4$ geometry with zero
free parameters.  The six predictions compiled here---a golden-ratio Bell
bound, a specific cosmic birefringence angle, a small tensor-to-scalar
ratio, a leptonic CP phase, golden-ratio gravitational-wave echoes, and
a proton lifetime---span quantum foundations, neutrino physics, cosmology,
gravitational-wave astronomy, and grand unification.  Each can be
confirmed or refuted by experiments expected to report within the next
five to ten years.  The framework puts itself at risk of falsification;
the tests are underway.

%----------------------------------------------------------------------
\section*{Data Availability}
%----------------------------------------------------------------------

All derivations, verification scripts (Python), formal proofs, and
simulation code are publicly available at
\url{https://github.com/grapheneaffiliate/e8-phi-constants}
under a CC-BY-4.0 license.

\begin{acknowledgments}
The author thanks the developers of Lean~4 for formal verification tools
and the LANL quantum noise group for publicly releasing ASE data.
\end{acknowledgments}

\begin{thebibliography}{20}

\bibitem{mcgirl2026gsm}
T.~McGirl,
``The Geometric Standard Model: E$_8 \times$ H$_4$ Unification of
Fundamental Constants,''
(2026), GitHub repository,
\url{https://github.com/grapheneaffiliate/e8-phi-constants}.

\bibitem{viazovska2016}
M.~S.~Viazovska,
``The sphere packing problem in dimension 8,''
\emph{Ann.\ Math.}\ \textbf{185}, 991 (2017).

\bibitem{mcgirl2026bell}
T.~McGirl,
``The Pentagonal Prism Bell Bound: A Golden-Ratio CHSH Inequality
from H$_4$ Coxeter Geometry,''
(2026), \url{https://github.com/grapheneaffiliate/e8-phi-constants}.

\bibitem{hensen2015}
B.~Hensen \emph{et al.},
``Loophole-free Bell inequality violation using electron spins
separated by 1.3 kilometres,''
\emph{Nature}\ \textbf{526}, 682 (2015).

\bibitem{codata2022}
E.~Tiesinga, P.~J.~Mohr, D.~B.~Newell, and B.~N.~Taylor,
``CODATA recommended values of the fundamental physical constants: 2022,''
\emph{Rev.\ Mod.\ Phys.}\ \textbf{95}, 025010 (2023).

\bibitem{pdg2024}
R.~L.~Workman \emph{et al.} (Particle Data Group),
``Review of Particle Physics,''
\emph{Phys.\ Rev.\ D}\ \textbf{110}, 030001 (2024).

\bibitem{minami2020}
Y.~Minami and E.~Komatsu,
``New extraction of the cosmic birefringence from the Planck 2018
polarization data,''
\emph{Phys.\ Rev.\ Lett.}\ \textbf{125}, 221301 (2020).

\bibitem{eskilt2022}
J.~R.~Eskilt,
``Frequency-dependent constraints on cosmic birefringence from the
LFI and HFI Planck Data Release 4,''
\emph{Astron.\ Astrophys.}\ \textbf{662}, A10 (2022).

\bibitem{planck2018}
N.~Aghanim \emph{et al.} (Planck Collaboration),
``Planck 2018 results. VI. Cosmological parameters,''
\emph{Astron.\ Astrophys.}\ \textbf{641}, A6 (2020).

\bibitem{bicepkeck2021}
P.~A.~R.~Ade \emph{et al.} (BICEP/Keck Collaboration),
``Improved constraints on primordial gravitational waves using
Planck, WMAP, and BICEP/Keck observations through the 2018 observing
season,''
\emph{Phys.\ Rev.\ Lett.}\ \textbf{127}, 151301 (2021).

\bibitem{abedi2017}
J.~Abedi, H.~Dykaar, and N.~Afshordi,
``Echoes from the abyss: Tentative evidence for Planck-scale structure
at black hole horizons,''
\emph{Phys.\ Rev.\ D}\ \textbf{96}, 082004 (2017).

\bibitem{superk2020}
A.~Takenaka \emph{et al.} (Super-Kamiokande Collaboration),
``Search for proton decay via $p \to e^+\pi^0$ and $p \to \mu^+\pi^0$
with an enlarged fiducial volume in Super-Kamiokande I--IV,''
\emph{Phys.\ Rev.\ D}\ \textbf{102}, 112011 (2020).

\bibitem{forbes2025}
I.~Nape, B.~Sephton, Y.-W.~Huang, A.~Vall\'{e}s, C.-W.~Qiu,
A.~Ambrosio, F.~Capasso, and A.~Forbes,
``Revealing the invariance of vectorial structured light in complex media,''
\emph{Nat.\ Commun.}\ \textbf{16}, 11305 (2025).

\end{thebibliography}

\end{document}
